% Abstract

\chapter*{Abstract}

In this manual, you get an introduction in typesetting a thesis with latex and
especially with the UniLeipzigBioinfThesis class.  The core is the use of the
class. However, there are also some hints and suggestions on what you should
do when you write your thesis. It starts with an introductory chapter.
Followed by three chapters that cover the main parts of a thesis: front
matter, main text and back matter.

It is designed to fulfill the \quo{Promotionsordung} from the \quo{Institut
für Informatik} at the \quo{Universität Leipzig}. The result is a class that
can be used without any further tweaking and create a good looking thesis that
contains all required parts. However, you can modify many of the design
decisions so that you can adapt it to your flavor.

Thus, the class helps you to directly start writing and not think too much
about the typesetting. I hope you enjoy this class.

% \todo{Starting here, the serious content ends and only blind text follows.}

% Here, I cite the biblatex website for test purposes~\citep{web:biblatex}.
% You should really read \citet{web:biblatex} as it contains valuable
% information.
%
% \Blindtext[8]

% End of frontmatter/abstract.tex
